\documentclass[11pt]{article}
\usepackage{enrico}
\usepackage{booktabs}
\usepackage{enumitem}
\usepackage[hidelinks]{hyperref}
\setlist{itemsep=2pt,topsep=3pt}
\numberwithin{theorem}{section}
\theoremstyle{definition}
\newtheorem{implementation}[theorem]{Implementation}
\newtheorem{algorithm}[theorem]{Algorithm}
\newtheorem{example}[theorem]{Example}
\theoremstyle{problemstyle}
\newtheorem{experiment}{Experiment}[section]
\nc{\ITT}{\mathrm{ITT}}
\nc{\LOO}{\mathrm{LOO}}
\nc{\ES}{\mathrm{ES}}
\nc{\SP}{\mathrm{SP}}
\nc{\sgn}{\operatorname{sign}}
\nc{\status}[1]{\hfill{\color{muted}\small[#1]}}
\definecolor{muted}{RGB}{105,112,122}

\title{Sandbox 1: organized populations of program-writing agents on Hanabi\\[2pt]\large Research notes}
\author{Enrico Yao-Bate}
\date{October 7, 2026}

\begin{document}
\maketitle
\tableofcontents

\section*{Preface}
These notes fix what Sandbox~1 is, in the order one needs to understand it: the goal and the claim we are making
(\S\ref{sec:goal}); the game, the bots, and the population process built on top of them (\S\ref{sec:game}); the causal
framework in which every quantity we report is defined (\S\ref{sec:causal}); the estimators and metrics and the rigor
protocol (\S\ref{sec:est}); the experiments, written as designs rather than results (\S\ref{sec:exp}); and the state of
the build (\S\ref{sec:build}). Each experiment carries a status tag, [not run], [pilot only] or [run]; no result
numbers appear in this version, since the experiments run so far were pilots used to fix designs, and their numbers
live in the ledger (\texttt{files/prereg/LEDGER.md}) and the status document. Definitions are given before they are
used; where a definition comes from a source, the source is cited by section (Wager, \emph{Causal Inference: A
Statistical Learning Approach}, draft of September 2026, cited as W \S; Sutton and Barto, \emph{Reinforcement
Learning}, second edition, cited as SB p.). Claims of novelty are phrased as ``not found'' with the nearest works named
in Appendix~\ref{app:novelty}.

\section{Goal, hypothesis, edge}\label{sec:goal}

The company is building the layer that lets a population of agents discover knowledge, share it, and keep building on
it over long horizons: who teaches whom, how a learner checks what it is taught, how credit and compute flow. Sandbox~1
is the controlled laboratory for that layer; Sandbox~2 (market making) points the same layer at the real domain.
Whatever we build here has to carry over there with only the evaluator swapped, which is the first constraint on
everything below.

\begin{ass}[The thesis, in testable form]\label{ass:thesis}
A long-running population of agents improves collectively because of (i) adaptive social dynamics: groups compete for
finite resources, so helping one's group is rewarded even when it costs the individual, and agents revise how they
cooperate and migrate between groups; and (ii) intelligence transfer: agents are credited for the improvement of those
they teach, discoveries from different lineages are combined, and learners verify what they are taught before adopting
it.
\end{ass}

\begin{rem}
Transfer alone is not interesting; many systems move knowledge between agents. The claims that distinguish us, stated
as the literature map states them, are five: the group as the unit that competes and reproduces, with teaching; causal
credit over a teaching graph with several teachers; controlled perturbation of the organization over long runs; a
formal model that predicts when organization beats scale; and designing against persuasive-but-wrong ideas and against
convergence. Rigor is itself part of the edge, since the adjacent literature runs single seeds, no counterfactuals and
no pre-registration. Of the five, the ones that are cheap enough to test before funding are credit, the formal model,
and robustness; the group-level and long-run claims need a factorial over organization rules and wait.
\end{rem}

\section{The game and the population}\label{sec:game}

\subsection{Hanabi}

Hanabi is a cooperative card game; we play it with two players, who form one team against the deck. Nobody in the game
competes with anybody, and no group exists inside a game. The goal is to build five fireworks, one per colour, by
playing cards onto each colour's stack in the order $1,2,3,4,5$; the score is the number of cards successfully played,
at most $25$. The twist is that each player holds their five cards facing away from themselves: a player sees the
partner's hand and never their own, so the only way to learn what one holds is from the partner's clues.

The deck has $50$ cards, five colours by ranks $1$ to $5$, with three $1$s, two each of $2,3,4$ and one $5$ per
colour. The table holds eight clue tokens and three life tokens. A turn is exactly one of three actions. A
\emph{clue} spends a token and names one colour or one rank, pointing at every card in the partner's hand that
matches. A \emph{discard} puts one of one's own cards face up, regains a token, and draws. A \emph{play} puts one of
one's own cards on the stacks; if it is the next rank of its colour it stays and scores, otherwise it is a bomb and
costs a life; either way one draws. The game ends when all stacks are complete, when three bombs have been played (our
engine scores this $0$), or when the deck runs out and each player has taken one last turn.

Eight tokens cannot describe $25$ cards literally, so good teams agree in advance on rules that make a clue mean more
than it says.

\begin{example}[A play clue]
The partner's newest card is a red $1$. We clue ``red'', touching only that card. Literally, the clue says that the
card is red. Under the convention ``a clue touching a single unknown card means play it'', the partner plays the card
without knowing its rank, and it scores. If the partner instead holds the convention ``a colour clue means save this
card'', they hold it forever. The same clue produces two outcomes depending only on whether the rule is shared, which
is the property that makes the game a laboratory for culture: a convention has no value alone, and whether two players
share one is measurable by letting them play.
\end{example}

The human community on \texttt{hanab.live} maintains a written rulebook of such conventions, in levels $1$ to $25$,
revised by pull request; expert pairs score about $23$ to $24$.

\subsection{Bots as policies}

\begin{definition}[Deal, policy]
A \emph{deal} is one shuffled order of the $50$ cards, indexed by a seed. A \emph{policy}, or \emph{bot}, is a
deterministic program $\pi$ mapping what one player can see (own hand knowledge, the partner's cards, the stacks, the
discards, the tokens, the last moves) to a legal move. A bot is reset at the start of every game and keeps no memory
across games.
\end{definition}

A game is thus a pair of bots and a deal. Nothing carries over between games, so evaluating a bot means playing it in
many independent games from fresh deals and averaging.

\begin{definition}[Score functional, self-play, cross-play]
For bots $\pi,\pi'$ and a finite set of deals $\D$, write
\[
\hat V_{\D}(\pi,\pi') := \frac{1}{|\D|}\sum_{d\in\D}\mathrm{score}(\pi,\pi';d),
\]
the mean score when $\pi$ and $\pi'$ play together over $\D$, with seats alternating by deal. $\hat V(\pi,\pi)$ is the
\emph{self-play} score of $\pi$; $\hat V(\pi,\pi')$ for $\pi\neq\pi'$ is the \emph{cross-play} score, the measure of
whether two bots share conventions. We write $V$ for the population version over the deal distribution.
\end{definition}

\begin{rem}
Self-play rewards internal consistency and says nothing about compatibility: two bots can each score $24$ with
themselves and $2.5$ with each other (Hu et al., \emph{Other-Play}, 2020), because each converged on private
signalling rules. This collapse is the quantitative face of ``independent groups develop incompatible cultures'', and
it is the null against which transfer claims are measured.
\end{rem}

\begin{definition}[Anchors]
An \emph{anchor} is a fixed reference bot with a known score, used so that scores are comparable across experiments
and time. Ours are: random play ($0$); the rule-based bots of Canaan et al.\ (AIIDE 2020), Piers ($16.99$ in two-player
self-play), IGGI ($15.99$) and Flawed ($0$, a deliberately bad rule list); and, optionally, a frozen reinforcement
learning policy.
\end{definition}

\begin{example}[Piers]
Piers encodes one human player's style as an ordered list of rules, the first applicable rule deciding the move: play a
card known to be safe; play the card most likely playable when a life can be spared; hint a partner's playable card;
hint useless cards when clue tokens are low; otherwise discard. Our engine reproduces its published $16.99$ over
$1{,}000$ deals move for move, which is how the engine adapter was validated. In the experiments below it is the
\emph{known-good teacher}: when a student adopts it, we know what the student should score.
\end{example}

\begin{rem}[Why this game]
Strategies are explicit rules, writeable as text and as code; a rule has value only when the partner shares it, so
knowledge must be transmitted to be worth anything; compatibility is measurable by cross-play; the score is absolute
with human and reinforcement-learning reference points; and a game simulates in about a millisecond, so we can play
millions. Games whose play requires a language model on every turn (social deduction games, Diplomacy) fail the last
property by three orders of magnitude and were rejected on that ground.
\end{rem}

\subsection{The population process}

We now describe the layer that sits on top of the game. The picture to keep in mind is a population of programmers,
each maintaining their own Hanabi bot, who can send one another their code and notes; the game is only the test bench.

\begin{definition}[Agent, incumbent, artifact]
An \emph{artifact} $a$ is a package of bot code together with a \emph{conventions document}, a plain-language
statement of the rules the code implements, with an identifier. An \emph{agent} is a process, driven by a language
model, that owns exactly one current artifact, its \emph{incumbent}, and can revise it, send it, and receive others'.
\end{definition}

\begin{definition}[Message, evidence]
A \emph{message} $m=(a,e)$ from a teacher $i$ to a student $j$ is an artifact together with \emph{evidence} $e$: the
scores our harness measured for $a$ (self-play, cross-play against anchors), signed by the harness. An agent cannot
fabricate evidence, only earn it.
\end{definition}

\begin{definition}[Verification, adoption]
On receiving $m=(a,e)$, the student may \emph{verify} $a$ by playing it against its own incumbent on a fixed set of
deals, and then \emph{adopt} it (replace the incumbent), \emph{merge} it (one language-model call combining the two),
or \emph{reject} it. Which of these happen, and under what rule, is an organization policy.
\end{definition}

\begin{definition}[Corpus, touch log]
The \emph{corpus} $\C_t$ of a group is the store of artifacts with evidence deposited by its members up to generation
$t$. A deposit is accepted only with harness evidence; a retrieval is recorded in the \emph{touch log} as (reader,
artifact, generation). The touch log is what makes credit computable later.
\end{definition}

\begin{definition}[Generation, lineage]
A \emph{generation} is one round in which every agent executes Algorithm~\ref{alg:step}. Every artifact records its
\emph{parents}: the agent's previous incumbent, together with any artifact adopted or merged. Following parents
backward gives an artifact's \emph{ancestry}; forward gives its \emph{descendants}, which include other agents' bots
that adopted it. The directed graph so formed is the \emph{lineage graph}.
\end{definition}

\begin{algorithm}[The agent step]\label{alg:step}
For one agent in one generation, in this order:
\begin{enumerate}
\item \emph{Receive} the messages the routing rule delivered to it.
\item \emph{Verify} each received artifact on the verification deals; adopt, merge or reject by the adoption rule.
\item \emph{Revise}: one language-model call that rewrites the incumbent, given the current code and conventions, a
summary of performance over the evaluation deals, the traces of the three worst games on the feedback deals, and
whatever was received and not withheld.
\item \emph{Evaluate} the new incumbent: self-play on the evaluation deals, cross-play against group members and
anchors.
\item \emph{Teach}: send the incumbent, with its evidence, to the agents the routing rule names.
\end{enumerate}
Each language-model call debits the agent's share of its group's budget; evaluation is engine time only.
\end{algorithm}

\begin{rem}[Two levels that never mix]
Inside a game two bots play sixty turns, deciding every move themselves, instantly, with no language model and no
other agent involved. The population never plays a game together, and nothing about it is decided by consensus or
vote: each agent decides what to do with its own bot, and the only interactions are messages, verification, adoption,
budget, and group membership. The wall-clock of a generation is dominated by the language-model calls, about a minute
per agent run one at a time, so twelve agents take about twelve minutes per generation and fifty generations take a
night; the games themselves take seconds.
\end{rem}

\begin{rem}[The right analogy]
The closest established frame is a genetic-programming population: artifacts are individuals, the language model is the
mutation operator, verification and adoption are selection, and teaching is horizontal transfer. What we vary is not
the operator but the organization around it: routing, verification, credit, allocation, migration.
\end{rem}

\subsection{Groups}

\begin{definition}[Group, group score]\label{def:group}
A \emph{group} is a subset of agents whose messages are routed mostly to one another and who share a corpus. Its
\emph{group score} is the expected score when two of its members, drawn at random, play together, i.e.\ the mean
within-group cross-play. Its \emph{compatibility} with another group is the mean between-group cross-play.
\end{definition}

\begin{rem}
This makes ``culture'' operational. A group whose members share conventions has a high group score; a group of
individually strong but mutually incompatible bots has a low one. Two isolated groups drift to different conventions,
so between-group cross-play falls toward the self-play collapse of the remark above; teaching across groups and
migration (an agent moving to another group, bringing its artifact) are the mechanisms that keep groups compatible,
and the lineage graph records which one did the work. Groups compete for the budget of language-model calls under an
allocation rule applied to group scores.
\end{rem}

\subsection{What varies}

Everything above is fixed across experiments. What an experiment varies is a tuple of organization policies, each a
small pluggable rule: topology (isolated, full, islands with migration); routing (none, broadcast within group, best to
all, random); delivery (deterministic, or Bernoulli with known probability); verification (none, self-play on $N$
deals, sequential); adoption (replace if better, merge, provisional with revert); credit (the estimators of
\S\ref{sec:est}); allocation (uniform, proportional, Nash-relative); migration; environment change (a rule variant at a
given generation); and the horizon, the number of generations for a fixed total budget.

\section{The causal framework}\label{sec:causal}

Every quantity we report is a causal effect of something the organization controls, on something the engine measures.
This section fixes the units, treatments and outcomes, and the assumptions under which the effects are identified.

\subsection{Potential outcomes for one student}

Consider one student with incumbent $\pi_0$, receiving $k$ messages indexed by $j\in[k]$. The organization decides
which messages are delivered; the student's update then runs. The obstacle to calling the result a causal effect of
delivery is that the update involves a language model, whose randomness would make the counterfactual unobservable.
The following assumption removes that obstacle, and it is the reason the local model is run with a fixed sampling seed.

\begin{ass}[Deterministic update]\label{ass:det}
The student's update $\pi' = U(\pi_0, S, \xi)$, where $S\subseteq[k]$ is the delivered set and $\xi$ the sampling
randomness of the language model, is a deterministic function of $S$ once $\xi$ is fixed. The engine is deterministic
given the deal.
\end{ass}

\begin{definition}[Delivery set, outcome]
Write $W\in\{0,1\}^k$ for the delivery indicators and $S=\{j: W_j=1\}$. The \emph{outcome} is the student's paired
improvement on the evaluation deals,
\[
Y(S) := \hat V_{\D_{\mathrm{eval}}}\big(\pi'(S),\pi'(S)\big) - \hat V_{\D_{\mathrm{eval}}}(\pi_0,\pi_0),
\]
with $\D_{\mathrm{eval}}$ disjoint from the deals shown to the model (\S\ref{sec:rigor}). A \emph{replicate} $r$ is
a change of experiment seed, which changes the feedback deals the model sees but not $\D_{\mathrm{eval}}$; we write
$Y_r(S)$.
\end{definition}

Under Assumption~\ref{ass:det}, $\{Y_r(S): S\subseteq[k]\}$ are the $2^k$ potential outcomes of replicate $r$, in the
sense of W \S1, and all of them can be observed by replaying. This is stronger than anything the causal-inference
literature assumes, and it is what makes an exact oracle possible (\S\ref{sec:est}).

\begin{definition}[Intent-to-treat effect of a message]\label{def:itt}
The intent-to-treat effect of message $j$ in replicate $r$ is
\[
\tau_{j,r} := \frac{1}{2^{k-1}}\sum_{S\subseteq[k]\setminus\{j\}}\big[\,Y_r(S\cup\{j\}) - Y_r(S)\,\big],
\]
the average, over the other messages delivered independently with probability $\tfrac12$ each, of the gain from
delivering $j$.
\end{definition}

The word ``intent'' is deliberate. Delivery is the treatment, and whatever the student does with the message, verify,
reject, adopt, or merely read, is inside the effect. This is the quantity the organization can act on, since delivery
is the only thing it controls, and it is the estimand that the pre-registered decision rules use.

\begin{prop}[Identification under randomized delivery]\label{prop:ident}
Let the organization deliver each message independently with $\P(W_j=1)=\tfrac12$. Then for every $j$,
\[
\E[Y\mid W_j=1] - \E[Y\mid W_j=0] = \tau_j ,
\]
where the expectation is over $W$ for fixed $r$ (and over $r$ as well if replicates are pooled).
\end{prop}
\begin{proof}
Conditional on $W_j=1$, the other indicators are still independent Bernoulli($\tfrac12$), so $\E[Y\mid W_j=1] =
2^{-(k-1)}\sum_{S\not\ni j}Y(S\cup\{j\})$, and similarly $\E[Y\mid W_j=0] = 2^{-(k-1)}\sum_{S\not\ni j}Y(S)$.
Subtracting gives the definition.
\end{proof}

\begin{rem}[What the sandbox gives for free]
Two things that are usually the hard part of causal inference are satisfied by construction here. First, the
propensity $\P(W_j=1)$ is known exactly, because we set it; W Corollary~3.3 (p.~38) says that with the true propensity
known, the augmented inverse-propensity estimator is efficient with any consistent regression model, so no nuisance
estimation is needed. Second, the deals are shared across every arm, so outcomes are paired, and W Theorem~1.3
(p.~14) says adjustment for a pre-treatment covariate is never worse than the raw difference in means; blocking on the
deal is therefore done by default. What is not given for free is the absence of interference between agents, taken up
in \S\ref{sec:interf}.
\end{rem}

\subsection{Credit as a game-theoretic quantity}

A second way to ask ``how much did message $j$ contribute'' is to divide the total gain fairly among the messages.

\begin{definition}[Shapley credit]\label{def:shapley}
With $v_r(S):=Y_r(S)$ as the value of a coalition of delivered messages,
\[
\phi_{j,r} := \sum_{S\subseteq[k]\setminus\{j\}} \frac{|S|!\,(k-|S|-1)!}{k!}\,\big[\,v_r(S\cup\{j\}) - v_r(S)\,\big].
\]
\end{definition}

\begin{prop}\label{prop:banzhaf}
$\tau_{j,r}$ is the Banzhaf index of $j$ in the game $v_r$, and $\sum_j\phi_{j,r} = Y_r([k]) - Y_r(\emptyset)$. If
$v_r$ is additive, $v_r(S)=\sum_{j\in S}a_j$, then $\tau_{j,r}=\phi_{j,r}=a_j$. In general they differ only through
the interaction terms of $v_r$.
\end{prop}
\begin{proof}
The Banzhaf index weights every marginal $v(S\cup\{j\})-v(S)$ equally by $2^{-(k-1)}$, which is
Definition~\ref{def:itt}; the Shapley weights sum to one and the efficiency identity is the standard telescoping over
a random ordering. Under additivity every marginal equals $a_j$, so any convex combination of them does.
\end{proof}

\begin{rem}[Which to use when]
The two answer different questions. $\tau$ answers ``what happens if I deliver $j$ under my delivery rule'', and it
changes if the rule changes, which is a feature when the rule is the object of study. $\phi$ answers ``how to split the
gain'', and its efficiency property matters as soon as credit is converted into budget: one cannot hand out more
credit than there was gain, and $\tau$ effects do not sum to anything meaningful once messages interact. Both come
from the same $2^k$ replays, so both are reported; the decision rules use $\tau$, because the organization's question
is causal rather than divisional. $\phi$ is also the formalism of the nearest published work on credit in multi-agent
language-model systems (Appendix~\ref{app:novelty}), which keeps our numbers comparable.
\end{rem}

\begin{definition}[Effect among adopters]
Delivery is an instrument for adoption: a message cannot be adopted unless delivered (one-sided noncompliance), so the
effect of $j$ among students who adopt it is identified as $\tau_{j}/\P(\text{adopt } j\mid j \text{ delivered})$
(W \S10). It is reported for interpretation and is undefined when no replicate adopts; it is not used in decision
rules, because conditioning on adoption directly would select on a post-treatment choice.
\end{definition}

\begin{rem}[Two lessons from the pilot that shape the designs]
First, a re-sent copy of the student's own bot is not a null message: any delivered text changes the revision, so a
``placebo'' with zero effect does not exist, and planted truths are only known in isolation. The oracle therefore
defines truth. Second, verification as implemented gates the adoption of code, not the reading of prose: a message
whose code is rejected can still move the student through the text the revision call sees. The reaction rules of
\S\ref{sec:exp} exist because of this.
\end{rem}

\subsection{Interference between agents}\label{sec:interf}

With several students in a group, one agent's treatment changes what is available to the others: a teacher's artifact
is the output of someone else's recent generation, and a verification or budget decision for one agent changes the
pool the rest draw from. The no-interference assumption behind \S3.1 fails across agents, and W \S11--12 give the
replacement.

\begin{ass}[Exposure mapping]
For agent $i$ there is a map $H_i:\{0,1\}^n\to\mathcal H$ such that $Y_i(w)=Y_i(w')$ whenever $H_i(w)=H_i(w')$ (W
Assumption~11.1, p.~146). Under network interference, $H_i(w)$ is the restriction of $w$ to $\{i\}\cup N_i$, the agent
and the set of agents whose artifacts or budget decisions can reach it in that generation (W Definition~11.1, p.~147).
\end{ass}

\begin{definition}[Randomization dependency graph]
$G_{ij}=1$ if $(\{i\}\cup N_i)\cap(\{j\}\cup N_j)\neq\emptyset$ (W eq.~12.18, p.~163). Outcomes of agents not joined
in $G$ are independent under the randomization.
\end{definition}

\begin{rem}
Three consequences for the group-level experiments. The inverse-probability exposure estimator is unbiased under
Bernoulli randomization with known exposure probabilities (W Theorem~12.1, p.~157), and ours are known. Its variance
should be the finite-population one (W Theorem~12.2, p.~159), which needs no story about a superpopulation of agent
cohorts; the network version (W eq.~12.21) is made conservative by projecting $G$ onto the positive semidefinite cone
(W Corollary~12.6, p.~167). Before assuming interference away, one tests for it: the permutation test for the sharp
null (W Theorem~11.1, p.~149) and the composite test for ``no spillovers'' (W Theorem~11.2). Inference remains valid
while the reach of any one agent grows slower than $n^{1/4}$ (W Theorem~12.7, p.~168), which caps how many students a
single teacher may touch in one generation relative to the population size, a design constraint on routing.
\end{rem}

\subsection{Evaluating an organization rule from logs}

A routing or verification rule is a policy over deliveries across generations. Running every candidate rule is
expensive; evaluating a candidate from logs collected under another rule is off-policy evaluation.

\begin{definition}[Behaviour and target rule]
Let $b$ be the rule that generated the logs and $\pi$ the rule to be valued, each assigning probabilities to deliveries
given the state. Coverage requires $\pi(a\mid s)>0 \Rightarrow b(a\mid s)>0$ (SB p.~103).
\end{definition}

\begin{prop}[Weighted importance sampling]
With $\rho_{t:T-1}=\prod_{u=t}^{T-1}\pi(A_u\mid S_u)/b(A_u\mid S_u)$ (SB eq.~5.3), the weighted estimator
$\sum_t \rho_t G_t/\sum_t\rho_t$ (SB eq.~5.6) is consistent under coverage and has bounded variance even when the
ratios have infinite variance; the ordinary estimator (SB eq.~5.5) is unbiased but can have unbounded variance.
\end{prop}

\begin{rem}
Over many generations the sequential weights compound multiplicatively (W eq.~14.12--14.14, pp.~195--196), the curse
of horizon, and the estimator becomes useless after a handful of generations. If the student state mixes, i.e.\
artifacts do not retain unbounded memory of old messages, the process is a Markov decision process and the doubly
robust stationary estimator of W \S15 (eq.~15.9, p.~208) avoids the compounding. If a rule is itself chosen
adaptively across generations, sample means are biased (W pp.~76--77); the adaptively weighted estimator (W eq.~6.12)
or a switchback design (W Definition~15.2, bias at most $4M\lambda(1+t_0)$, Theorem~15.5) restores valid inference.
These are the tools for the long-run experiments, not for the single-student ones.
\end{rem}

\subsection{Credit over time, and teacher choice}

\begin{rem}[Eligibility traces as lineage credit]
Crediting a teacher for a student's improvement several generations downstream is Minsky's credit-assignment problem
(SB p.~17) with messages in place of decisions. The backward view of TD($\lambda$) (SB eqs.~12.5--12.7, pp.~292--294)
assigns the current error backward along the trajectory with weight decaying by $\gamma\lambda$ per step; a trace per
(teacher, student) pair along the lineage graph is the direct analogue, with $\lambda=1$ recovering plain Monte Carlo
credit when the bootstrapped version is not trusted. The book's traces run along one trajectory; with several teachers
acting on one student the per-pair trace is the minimal extension, and this is a design choice we make rather than a
result we inherit.
\end{rem}

\begin{rem}[Which teacher to listen to]
Choosing teachers under uncertainty is a nonstationary bandit, since teachers improve. The constant-step-size update
$Q_{n+1}=Q_n+\alpha(R_n-Q_n)$ (SB eq.~2.5, p.~32) is appropriate rather than the sample average, and the book's own
caveat applies: upper-confidence selection assumes uncertainty shrinks with the pull count, which fails under
nonstationarity (SB p.~36). Any baseline that does not depend on the action leaves the expected update unchanged and
reduces its variance (SB p.~329). This mechanism needs many generations to matter and is deferred.
\end{rem}

\section{Estimators, metrics, and the rigor protocol}\label{sec:est}

\subsection{The oracle}

\begin{definition}[Oracle]
For replicate $r$, run the student's update for every $S\subseteq[k]$ under the same seed, obtaining all $2^k$ values
$Y_r(S)$; compute $\tau_{j,r}$ and $\phi_{j,r}$ from Definitions~\ref{def:itt} and \ref{def:shapley}.
\end{definition}

\begin{rem}
The oracle costs $2^k$ language-model calls per replicate: $8$ at $k=3$, $128$ at $k=7$. It is affordable for small
$k$ and becomes the problem at large $k$, which is why cheaper estimators are tested against it. Its existence depends
entirely on Assumption~\ref{ass:det}; with a model that cannot be seeded one would replace each $Y_r(S)$ by an average
over repeated calls, multiplying the cost.
\end{rem}

\subsection{Estimators}

Each estimator uses a subset of the $2^k$ replays, and the question is how much accuracy is lost.

\begin{definition}[Leave-one-out]
$\hat\tau^{\LOO}_{j,r} := Y_r([k]) - Y_r([k]\setminus\{j\})$, using $k+1$ replays.
\end{definition}

\begin{prop}
$\hat\tau^{\LOO}_{j,r}=\tau_{j,r}$ when $v_r$ is additive. In general it equals the single marginal of $j$ with all other
messages present, one of the $2^{k-1}$ terms averaged in Definition~\ref{def:itt}.
\end{prop}

\begin{rem}
This is the ablation of SHARP (arXiv 2602.08335), a 2026 credit scheme for multi-agent language-model systems that,
despite its name, masks one agent and replays; we generalize its binary accuracy to our continuous outcome. It is the
cheapest replay-based estimator with a precedent, and it is biased exactly when messages interact: the margin of a good
artifact is smaller when a persuasive bad one is also in the prompt. The opposite single context, leave-one-in,
$Y_r(\{j\})-Y_r(\emptyset)$, costs the same and has the mirror-image bias.
\end{rem}

\begin{definition}[Equal split]
Deliver everything once; let $A_r$ be the set of adopted messages. $\hat\tau^{\ES}_{j,r} := Y_r([k])\,\mathbf
1\{j\in A_r\}/|A_r|$, one replay.
\end{definition}

\begin{rem}
This is the accounting an organization does with no replays at all, and it is what delivery-based credit schemes
amount to. It assigns nothing to rejected messages, so it cannot see the effect of text that was read but not adopted,
nothing to differences in magnitude among adopted ones, and is undefined when nothing is adopted. It is the zero-cost
baseline every other estimator must beat.
\end{rem}

\begin{definition}[Randomized delivery with regression]\label{def:ridge}
Let the organization deliver each message by an independent coin flip $W_j\sim\Bern(\tfrac12)$ and log every event
$(W,Y)$. Pooling events over replicates $r$ and draws $i$, fit
\[
(\hat a,\hat\beta) := \arg\min_{a,\beta}\;\sum_{r,i}\Big(Y_r(W^{(i)}_r) - a - \textstyle\sum_{j}\beta_jW^{(i)}_{r,j}\Big)^2
+ \lambda\|\beta\|_2^2,
\qquad \hat\tau^{\mathrm{ridge}}_j := \hat\beta_j .
\]
\end{definition}

The model reads: the student's gain is a baseline $a$, what it gets from its own revision with nothing delivered, plus
a contribution $\beta_j$ for each message that arrived. The coefficients are on the coin flips because the coin flips
are the treatments. This is the regression form of a randomized experiment with several simultaneous treatments
(W \S1: the difference in means is the regression on the treatment indicator), and the datamodels estimator of Ilyas
et al.\ (2022) in the credit setting.

\begin{prop}
With $\lambda=0$ and independent Bernoulli($\tfrac12$) deliveries, $\hat\beta_j\to\tau_j$ as the number of events
grows, where $\tau_j$ is the replicate-averaged effect $\frac1R\sum_r\tau_{j,r}$.
\end{prop}
\begin{proof}
The design is orthogonal, so the least-squares coefficient on $W_j$ converges to $\E[Y\mid W_j=1]-\E[Y\mid W_j=0]$,
which is $\tau_j$ by Proposition~\ref{prop:ident}; interactions are averaged over the other coin flips, not assumed
away.
\end{proof}

\begin{rem}[Why ridge, and what it does and does not estimate]
The penalty $\lambda\|\beta\|^2$ shrinks the coefficients toward zero, trading a little bias for variance when events
are few relative to messages; at $k=3$ with hundreds of events it is immaterial, and it matters in the regime the
estimator exists for, many messages and one event per student-generation. Alternatives: plain least squares; the lasso
when most messages do nothing; interaction terms $W_jW_l$ to recover the pairwise structure that $\phi$ weights;
replicate fixed effects to absorb the large between-student baselines. The estimator costs no replays beyond the
deliveries the organization randomizes anyway, which is its reason to exist, and it estimates the \emph{average}
credit across students and generations, not any one student's. That distinction is the content of
Proposition~\ref{prop:floor}.
\end{rem}

\begin{definition}[Singles plus pairs]
Replay the empty set, each singleton and each pair, $1+k+\binom k2$ replays, and set
\[
\alpha_{j,r} := Y_r(\{j\})-Y_r(\emptyset),\qquad
\gamma_{jl,r} := Y_r(\{j,l\})-Y_r(\{j\})-Y_r(\{l\})+Y_r(\emptyset),
\]
with $\hat\tau^{\SP}_{j,r}$ the average of Definition~\ref{def:itt} in which any missing $Y_r(S)$ with $|S|\ge3$ is
imputed from $\alpha$ and $\gamma$ assuming no higher-order interaction.
\end{definition}

\begin{rem}
The replay pattern is that of Li et al.\ (2026, arXiv 2608.29228), who search singletons and pairs for the minimal
message sets that flip a failed run. It is the cheapest estimator that sees pairwise interactions; at $k=3$ its seven
replays are no saving over the oracle's eight, and it pays off at $k=7$ ($29$ against $128$). Designed delivery
patterns, Plackett--Burman arrays in place of coin flips, serve the same purpose for the regression estimator and are
tested in the same arm.
\end{rem}

\subsection{Metrics}

\begin{definition}[Per-replicate accuracy]
For an estimator $\hat\tau$ and replicate $r$,
\[
\mathrm{RMSE}_r := \Big(\tfrac1k\sum_j(\hat\tau_{j,r}-\tau_{j,r})^2\Big)^{1/2},\qquad
\rho_r := \mathrm{Spearman}\big((\hat\tau_{j,r})_j,(\tau_{j,r})_j\big),\qquad
e_r := \mathbf 1\{\sgn\hat\tau_{j^\ast,r}\neq\sgn\tau_{j^\ast,r}\},
\]
where $j^\ast$ is the planted harmful message. Each is aggregated over replicates by the interquartile mean, with a
$95\%$ interval from bootstrap resampling of replicates within stratum.
\end{definition}

\begin{rem}[Why these three]
Credit feeds budget allocation, so errors of magnitude cost compute; root mean square penalizes large misses more than
mean absolute error, which is the right sensitivity for that use. Routing decisions need only ordering, hence the rank
correlation; with three messages it takes only four values and is a coarse check, as Kendall's $\tau$ would be. The
sign error on the harmful message is the safety question: does credit at least say that harm was harmful. The
interquartile mean with bootstrap intervals is the recommendation of Agarwal et al.\ (2021) for small numbers of runs,
so that one bad seed cannot drive a verdict. A calibration slope, the regression of estimate on truth, would say whether
an estimator's error is a constant scale factor and should be added.
\end{rem}

\begin{prop}[Floor for pooled estimators]\label{prop:floor}
Let $\hat\tau_j$ be any estimator that returns the same value for every replicate. Then
\[
\E_r\big[(\hat\tau_j-\tau_{j,r})^2\big]\;\ge\;\Var_r(\tau_{j,r}),
\]
with equality when $\hat\tau_j=\E_r\tau_{j,r}$.
\end{prop}
\begin{proof}
Add and subtract $\bar\tau_j=\E_r\tau_{j,r}$ and expand; the cross term vanishes.
\end{proof}

\begin{rem}
The regression estimator of Definition~\ref{def:ridge} is pooled, so its per-replicate root mean square error is
bounded below by the between-replicate standard deviation of the truth, however many events it sees. A per-replicate
accuracy criterion therefore asks a pooled estimator a question it cannot answer, and the fair criterion for it is its
error against the replicate-averaged credit $\bar\tau_j$. Two targets follow: per-student credit, available only by
replay; and population-average credit, available from randomized delivery at no extra cost. Pre-registrations must say
which one an estimator is judged against.
\end{rem}

\subsection{Rigor protocol}\label{sec:rigor}

\begin{implementation}[Deal separation]
Every generation draws three disjoint sets of deals from the shared stream: \emph{feedback} deals, whose worst traces
the model sees; \emph{verification} deals, on which received artifacts are judged; \emph{evaluation} deals, on which
$Y$ is computed and which the model never sees. The generalization gap, evaluation score minus feedback-deal score, is
logged per artifact. Deals are independent draws from one fixed distribution, so a random split is the correct one;
the walk-forward discipline of time series is needed only across generations, where it takes the form of the frozen
ladder below.
\end{implementation}

\begin{implementation}[Pre-registration]
Before a run, one page fixes: the question; the setting; the estimand; the known or planted truth where one exists;
the design (what is paired, randomized, blocked); the metrics; the decision rules with thresholds; the minimum effect
worth detecting; a pilot of five seeds to estimate variance, then the number of seeds by a power calculation; the
stopping rule; what would falsify the approach; compute and cost. Amendments after the pilot are logged with a
timestamp, never silently edited. The ledger (\texttt{files/prereg/LEDGER.md}) keeps every pre-registration and its
outcome, negatives included.
\end{implementation}

\begin{implementation}[Null populations, evidence ladder, multiplicity]
The whole pipeline is run with a null stub that only adds noise, giving the null distribution of every reported metric;
figures carry null bands. A result is accepted only when the pre-registered effect's interval excludes zero at the
declared minimum effect on two local models (the second as a blocking factor). The Benjamini--Hochberg procedure at
$q=0.1$ is applied across the confirmatory tests with $p$-values in the ledger; threshold rules are acceptance
criteria, not tests, and exploratory contrasts are reported with their $p$-values outside the family. Every curve is
plotted against tokens and seconds; local runs cost nothing but are timed.
\end{implementation}

\begin{implementation}[Frozen ladder]
In population runs every generation's artifacts are kept, and the current population is cross-played against every
past generation and the anchor set, so improvement is measured against fixed opponents and partners rather than
against a moving population; the league Nash assigning small weight to past generations is the non-cycling check
(AlphaStar, Vinyals et al.\ 2019).
\end{implementation}

\section{Experiments}\label{sec:exp}

Each experiment states what it is for, what it is, why it is worth doing now, how it is done, its decision rule and
its cost. The single-student designs share one driver; the population designs share another. Status tags are [not
run], [pilot only] (a five-seed pilot was run to fix the design; its numbers are in the ledger and are not results),
and [run].

\begin{experiment}[C1: which credit estimator to trust]\status{pilot only; full run on one model block}
\emph{For.} Causal credit over a teaching graph, the second edge claim.
\emph{What.} A student with incumbent $\pi_0$ receives $k=3$ messages carrying Piers, Flawed with persuasive prose and
honest low evidence, and a copy of its own bot. The oracle replays all eight delivery subsets; the estimators of
\S\ref{sec:est} are scored against it.
\emph{Why.} Credit for students' improvement is what makes teaching in an agent's interest, and the nearest published
schemes have no empirical validation or no code; an estimator we trust every generation is a prerequisite for any
organizational experiment.
\emph{How.} Two strata, verification on (self-play on $200$ deals, replace if better) and off; replicates $R$ from a
five-seed pilot by the power rule at a one-point minimum effect on the primary contrast; a second arm with $k=7$ (four
null teachers added) for singles-plus-pairs and Plackett--Burman against an exhaustive oracle on five seeds; a
replication block on the second model.
\emph{Decision rule.} An estimator is accepted if, in both strata and both blocks, Spearman interquartile mean $\ge0.9$
with bootstrap lower bound $\ge0.8$, root mean square error $\le2$ points, and sign error on Flawed $\le5\%$; among
accepted, the cheapest. Primary contrast: error of the regression estimator minus error of leave-one-out, paired by
replicate.
\emph{Cost.} $2^3\times R\times2$ local calls.
\end{experiment}

\begin{experiment}[C1b: pooled estimators against the population-average credit]\status{not run}
\emph{For.} The same claim, with the criterion matched to what a pooled estimator estimates
(Proposition~\ref{prop:floor}).
\emph{What and how.} Same replays as C1; the regression estimator (coin flips, and Plackett--Burman patterns) is
judged by its error against $\bar\tau_j$ with a bootstrap interval over replicates, at one event per student-generation
to mimic a running organization.
\emph{Decision rule.} Accept if the error against $\bar\tau_j$ is at most one point with the interval below $1.5$, in
both blocks.
\end{experiment}

\begin{experiment}[C2: does withholding unverified text remove the influence of bad ideas]\status{pilot only}
\emph{For.} Robustness to persuasive-but-wrong ideas, the fifth edge claim.
\emph{What.} The setting of C1 with a \emph{quarantine} rule: a message reaches the revision prompt only after its
code passes verification. Under Assumption~\ref{ass:det} the effect of a withheld message is zero by construction, so
the quantity with empirical content is the effect of a rejected message's text \emph{without} quarantine: the contrast
between Flawed and the re-sent own bot under verification on, paired by replicate.
\emph{Decision rule.} Claim that rejected text influences the learner if that contrast's interquartile mean is at least
two points with the bootstrap interval excluding zero, in both blocks; claim quarantine is costless for good teachers if
the Piers effect under quarantine is within one point of its effect without.
\end{experiment}

\begin{experiment}[D1: breakdown point of verification, with reaction rules]\status{not run on a model}
\emph{For.} The fifth edge claim at the group level: how much contamination the organization tolerates.
\emph{What.} One group of $n$ students, $G$ generations, broadcast routing. Each delivered message is independently
replaced with probability $\varepsilon$ by a Flawed-family artifact with persuasive prose and honest evidence. Two
factors: $\varepsilon\in\{0,0.1,0.2,0.3,0.5\}$ and the \emph{reaction rule}, the rule by which a student treats what it
receives.
\emph{Why.} In robust statistics the breakdown point is the contamination an estimator tolerates; here the
organization is the estimator and bad messages are contamination. Byzantine-robustness results give fraction
thresholds for consensus rules and memory-poisoning defenses give attack rates; the breakdown curve for execution-based
verification in a learning population was not found (Appendix~\ref{app:novelty}).
\emph{How.} Sabotage is Bernoulli with known probability, so exposure probabilities are known (W Theorem~12.1).
Outcome: the group score at generation $G$ (Definition~\ref{def:group}), paired on deals. Breakdown point $\varepsilon^\ast$: the
smallest $\varepsilon$ at which the group score falls below $70\%$ of its uncontaminated value, interpolated between
grid points. Variance by the finite-population estimator with the projected dependency graph; permutation test of the
sharp null of no $\varepsilon$ effect within each rule. Pilot on the null stub, then five seeds per cell on the model
at reduced size, then $R$ by power.
\emph{Decision rule.} Claim that a rule raises the breakdown point if $\varepsilon^\ast(\text{rule}) -
\varepsilon^\ast(\text{read everything})\ge0.2$ with the interval excluding zero, in both blocks.
\end{experiment}

The reaction rules are defined here because they are the object D1 compares.

\begin{definition}[Reaction rules]
On receiving $m=(a,e)$, a student may: (i) \emph{read everything}: verify $a$, adopt by the adoption rule, and let the
revision call see $m$ regardless of the verdict (the baseline); (ii) \emph{quarantine}: let the revision see $m$ only
if $a$ passed verification; (iii) \emph{evidence-gated reading}: let the revision see $m$ only if the signed evidence
$e$ exceeds the student's own score; (iv) \emph{idea trial}: implement the conventions of $m$ in a sandboxed branch, test
the branch, and let the revision see $m$ only if the branch beats the incumbent; (v) \emph{trust weighting}: weight
$m$ by the teacher's credit history under a nonstationary bandit; (vi) \emph{critical social learning}: adopt
provisionally and revert if the next evaluation is worse (Enquist, Eriksson and Ghirlanda, 2007).
\end{definition}

\begin{rem}[Trade-offs among the rules]
Quarantine costs nothing and guarantees zero influence from unverified messages, but discards a good idea whose code
happens to fail, and delays reading until verification has run. Evidence-gated reading is free and uses only signed
numbers, but score is not relevance: a weaker teacher with one useful specific idea is ignored. The idea trial is
verification of ideas rather than of artifacts, the principled rule, at the price of one extra language-model call per
message, which doubles the cost of a teaching event. Trust weighting and critical social learning act over
generations rather than within one, so they need the population runs and are deferred; the former has a cold-start
and lock-in risk, the latter accepts transient damage and needs a paired test to revert reliably, which the sequential
test of B1 supplies.
\end{rem}

\begin{experiment}[B1: verification as a sequential test]\status{run, engine only}
\emph{For.} The cost side of verification.
\emph{What.} Each adoption decision is a paired test on per-deal differences $d_1,d_2,\dots$ between candidate and
incumbent, with mean $\delta$ and standard deviation $\sigma$. Wald's sequential probability ratio test continues
until the log-likelihood ratio leaves $[\log\tfrac{\beta}{1-\alpha},\log\tfrac{1-\beta}{\alpha}]$; its expected
sample size at error rates $\alpha,\beta$ is of order $2\sigma^2\log(1/\alpha)/\delta^2$.
\emph{How.} A bank of candidate--incumbent pairs with $\delta$ measured on $2{,}000$ deals; fixed $N\in\{50,100,200,400\}$
against the sequential test at matched error rates; false adoptions, misses near $\delta=1$, and deals used.
\emph{Decision rule.} Adopt the sequential test as the default verification rule if it matches the fixed-$N$ error rates
with at least $30\%$ fewer deals at $\delta=1$. Expected regret in score points should be reported alongside.
\end{experiment}

\begin{experiment}[A1: transmission fidelity and the accumulation prediction]\status{pilot only}
\emph{For.} The formal model, the fourth edge claim.
\emph{What.} Send a student a known-good artifact (Piers, self-play $v_T$) as code, as prose, or as both; the student
revises once; measure the copy. Copy loss $\alpha(\mu):=v_T-\E\,\hat V(S',S')$ and copy dispersion $\beta(\mu)$ per
medium $\mu$, and the sufficiency ratio $\rho(\mu):=\E\,\hat V(S',T)/v_T$, the cross-play of the copy with its
teacher relative to the teacher's self-play. Henrich's (2004) model predicts mean skill change per generation
$-\alpha+\beta(0.577+\ln N)$ when learners copy the best of $N$ exemplars with Gumbel copy error; the sign predicts
whether a small population improves, and a check population tests the sign.
\emph{Why.} Fidelity rather than population size decides accumulation in that theory; its parameterization and a
code-versus-prose comparison were not found for language-model agents.
\emph{How.} Thirty replicates per medium with the feedback deals varied (at temperature zero a sampling seed does not
vary the output); verification off; all copies scored on one common deal set. The pilot showed the copy distribution
to be a cluster at the teacher's level plus a tail of failed copies, so the Gumbel dispersion is the wrong summary: the
fidelity of a medium is to be reported as a failure probability plus a within-cluster loss, and the prediction derived
from that mixture before the check populations run.
\emph{Decision rule.} Medium claim if $\alpha(\text{prose})-\alpha(\text{code})\ge1$ with the interval excluding
zero; model-support claim if the predicted sign matches the observed sign in at least eight of ten population seeds
per medium, in both blocks.
\end{experiment}

\begin{experiment}[R1: recombination across lineages]\status{not run}
\emph{For.} The claim that teaching combines discoveries from different lineages.
\emph{What.} Plant two complementary partial conventions in two agents (one a good discard rule, one a good hint rule,
each alone worth little), and ask whether, and how quickly, the population produces an artifact that has both, under
isolated, full and island topologies. This is the two-path design of Derex and Boyd (2016) with program-writing agents.
\emph{How.} Outcome: generations until an artifact beats the better of the two parents by a pre-set margin on
evaluation deals, with the lineage graph showing both parents in its ancestry; null band from the null stub.
\end{experiment}

\begin{experiment}[H1: horizon at fixed budget]\status{not run; the lab's central hypothesis}
\emph{What.} The same total budget split as twenty generations with large per-generation budgets versus two hundred
with small ones, with and without teaching. If many small generations win only when teaching is on, long-running is
doing the work the thesis says it does.
\emph{How.} Frozen ladder; trajectory metrics (slope in the last third against the first third; changepoints;
retained innovations, artifacts that beat the archive and still have descendants $k$ generations later).
\end{experiment}

\begin{rem}[Deferred]
Groups competing for compute under allocation rules, migration rates, selection rules such as clade-level credit, and
learned routing are implemented as frozen, tested policies and are not run until the single-student and small-group
results exist; a factorial over them is a post-funding experiment.
\end{rem}

\section{Build and status}\label{sec:build}

\begin{implementation}[Engine and anchors]
The Hanabi Learning Environment is vendored at a pinned commit and built from source; bots see a fixed observation
schema and act through a sandbox (import whitelist, guarded attribute access, timeouts, a watchdog) so that a bot
cannot read its own cards or the file system. The adapter reproduces the published anchor scores move for move. Games
export to the \texttt{hanab.live} replay format, so any game can be watched.
\end{implementation}

\begin{implementation}[Local model and determinism]
Agents run on an open-weight mixture-of-experts model (Qwen3.6-35B-A3B, 4-bit) served by MLX on the development
machine, with sampling seed and temperature zero: identical requests return identical text across sessions, which is
Assumption~\ref{ass:det} in practice. A second model (Gemma~4 26B-A4B) is the replication block. No paid model is in
the loop. Calls run one at a time, since batching on the server breaks reproducibility.
\end{implementation}

\begin{implementation}[Harness]
Artifacts are content-addressed with harness-signed evidence; every language-model call is recorded and replayable;
runs checkpoint atomically and resume byte-identically; configurations are frozen by digest; the single-student driver
implements the oracle and every estimator of \S\ref{sec:est}; the population driver implements the policies of \S\ref{sec:game}.5
with the group-level ones frozen. The test suite exercises determinism, the sandbox, the estimators on synthetic truths,
the credit identities, and crash-resume.
\end{implementation}

\section{Next steps}
\begin{enumerate}
\item Pre-register D1 with the reaction-rule factor and C1b; run their pilots on the model.
\item Replicate C1 on the second model, so that an estimator verdict can be accepted rather than only rejected.
\item Re-register C2 around the rejected-text contrast; re-register A1 around the mixture summary of fidelity.
\item R1 and H1 on a small population once the single-student results are in.
\end{enumerate}

\appendix
\section{Novelty notes}\label{app:novelty}
From targeted searches of arXiv on 2026-10-06; venues were not checked separately. Phrasing in any public document
follows these.
\begin{itemize}
\item \emph{Credit.} Message-level credit in multi-agent language-model systems is an active area: Shapley with process
rewards (Yang et al.\ 2025, arXiv 2511.10687, conceptual); counterfactual replay for failure localization (Li et al.\
2026, arXiv 2608.29228); Aumann--Shapley attribution at scale (Tang et al.\ 2026, arXiv 2605.11404); SHARP (arXiv
2602.08335). Not found: a randomized-delivery regression checked against exact seeded replay as a teacher's causal
effect on a paired-deal student.
\item \emph{Breakdown point.} Fraction thresholds for consensus under corrupted agents (Liu et al.\ 2026, arXiv
2604.17139; Lee et al.\ 2026, arXiv 2605.09076); memory-poisoning gates with attack rates (MAPLE-Guard, arXiv
2608.00426). Not found: a breakdown curve for execution-based verification in a learning population.
\item \emph{Fidelity.} Transmission-chain studies of language models (Perez et al.\ 2024, arXiv 2407.04503 and
2403.08882); generational cooperation (Vallinder and Hughes 2024, arXiv 2412.10270); peer-verified accumulation (POLIS,
arXiv 2507.21166). Not found: Henrich's parameterization or a code-versus-prose comparison for agents.
\end{itemize}

\section{Glossary}
Deal: one shuffled deck. Policy, bot: a deterministic program from observation to move. Self-play, cross-play: score of
a bot with itself, with another. Anchor: a fixed reference bot. Artifact: code plus conventions document. Agent: a
language-model-driven owner of one artifact. Message, evidence: an artifact with harness-signed scores. Verification,
adoption: testing a received artifact; replacing one's own. Corpus, touch log: a group's store of artifacts; the
record of who read what. Generation: one round of Algorithm~\ref{alg:step}. Lineage graph: parents and descendants of
artifacts. Group, group score: a subset with shared routing and corpus; mean within-group cross-play. Replicate: a change
of experiment seed. Intent-to-treat effect, Shapley credit: Definitions~\ref{def:itt} and \ref{def:shapley}.
\end{document}
